\documentclass[]{beamer}
% \usetheme{metropolis} 
\usepackage{amsmath, amsfonts, amssymb}

\title{Journal Club}
\subtitle{Schulman, John, et al. "Proximal policy optimization algorithms." arXiv preprint arXiv:1707.06347 (2017).}
\author{Jongkook Choi}
\date{2026-04-29}

\begin{document}

\frame{\titlepage}

\begin{frame}{Introduction}

    \begin{itemize}
        \item Schulman, John, et al. "Proximal policy optimization algorithms." arXiv preprint arXiv:1707.06347 (2017).
        \item Cited by 38489
        \item[] 
        \item John Schulman : Co-founder of OpenAI.
    \end{itemize}

\end{frame}

\begin{frame}{Policy Gradient}
    \textbf{Conceptual Framework}
    \begin{itemize}
        \item \textbf{Objective:} Optimal control of stochastic trajectories by finding a policy $\pi_{\theta}$ that maximizes the expected cumulative reward $J(\theta) = \mathbb{E}_{\tau \sim \pi_{\theta}}[\sum r_t]$.
        \item \textbf{From SDEs to RL:} Applying the Markov assumption to stochastic processes allows us to parameterize the control law as a stochastic policy $\pi_{\theta}(a_t | s_t)$.
    \end{itemize}

    \textbf{The Policy Gradient Theorem}
    The gradient of the objective is estimated using the log-derivative trick (likelihood ratio):
    $$ \nabla_{\theta} J(\theta) = \mathbb{E}_{\tau \sim \pi_{\theta}} \left[ \sum_{h=0}^{H-1} \nabla_{\theta} \ln \pi_{\theta}(a_h | s_h) R_h \right] $$

\end{frame}

\begin{frame}{}
    \textbf{Supervised Learning}
    \begin{itemize}
        \item \textbf{Data-Driven:} Track individual "particles" (data points) to sample sequences from a fixed dataset.
        \item \textbf{Objective:} Train the Transformer using Maximum Likelihood Estimation (MLE) until prediction accuracy is maximized.
    \end{itemize}
    
    \vfill
    
    \textbf{Monte Carlo Approach (Policy Gradient/PPO)}
    \begin{itemize}
        \item \textbf{Feynman-Kac:} The expected outcome of a stochastic process (SDE) can be estimated by sampling trajectories dictated by a "drift" term (PDE).
        \item \textbf{Objective:} Train the Transformer to optimize this drift (the policy), shifting the distribution toward high-reward trajectories. (so-called delta hedge)
    \end{itemize}
\end{frame}

% \begin{frame}{Importance Sampling}

%     Monte Carlo method requires many sampling when the variance of the distribution can be high. (e.g. extreme reward, rare actions)

%     To reduce it, we can update the policy much when we got the rare observation. 

%     Measure theoretically speaking, 
%     $$E_P[X] = \int X p(x) dx = \int X dP = \int X dP/dQ dQ = E_Q[X dP/dQ]$$
    
%     dP/dQ here is the importance of the sampled step.
% \end{frame}

\begin{frame}{Importance Sampling: Change of Measure}
    \textbf{Motivation: Rare Events \& High Variance}
    \begin{itemize}
        \item Monte Carlo estimation suffers from high variance when target events are rare or the distribution has heavy tails.
    \end{itemize}
    ~\\
    \textbf{Change of Measure}
    To estimate $\mathbb{E}_P[f(X)]$ using samples from $Q$:
    $$ \mathbb{E}_P[f(X)] = \int f(x) p(x) dx = \int f(x) \frac{p(x)}{q(x)} q(x) dx = \mathbb{E}_Q \left[ f(X) \frac{dP}{dQ} \right] $$

    $\rightarrow$ Efficient data usage is possible when P is close enough to the Q

\end{frame}

\begin{frame}
    \textbf{Application in PPO}
    \begin{itemize}
        \item \textbf{Target ($P$):} The new policy $\pi_{\theta}$ we want to optimize.
        \item \textbf{Proposal ($Q$):} The old policy $\pi_{\theta_{old}}$ used for data collection.
        \item \textbf{The Ratio ($dP/dQ$):} The importance weight $r_t(\theta) = \frac{\pi_{\theta}(a_t|s_t)}{\pi_{\theta_{old}}(a_t|s_t)}$.
    \end{itemize}

    \textbf{Example}
    \begin{itemize}
        \item Calculate importance under Markovian assumption $\frac{dP}{dQ} = \frac{\pi_{\theta}(a_t|s_t)}{\pi_{\theta_{old}}(a_t|s_t)}$. (You can calculate it from the logits)
        \begin{itemize}
            \item You could add transformer context as state, but there's high chance of you facing the curse of dimensionaliy.
        \end{itemize}
        \item Sample trajectories using $\pi_{old}$
        \item Multiply weight using the importance, 
        \item Estimate the expected gradient using the law of large numbers  
    \end{itemize}
\end{frame}

\begin{frame}{Proximal policy optimization}
    "Proximal" : Keeping the new policy within a specific "proximity" or neighborhood of the old policy
\\~\\
    \textbf{PPO with Clipping and the Pessimistic Bound}\\
    For stability, PPO defines a pessimistic lower bound on the objective
    $$ L^{CLIP}(\theta) = \mathbb{E}_t \left[ \min \left( r_t(\theta) (R_t - B_t), \text{clip}(r_t(\theta), 1-\epsilon, 1+\epsilon) (R_t-B_t) \right) \right] $$
    \begin{itemize}
        \item Ignores improvement by extreme choice (gradient becomes zero when reward is greater than baseline and the importance is high) 
        \item Ignores bad luck when the choice is normal (gradient becomes zero when reward is less than baseline and the importance is high) 
    \end{itemize}
\end{frame}

\end{document}